\subsection{DES}
\subsubsection{DES简介}
数据加密标准（Data Encryption Standard, DES）是由美国国家标准局（National Institute of Standards and Technology, NIST）在1977年采纳作为“非密级”应用的分组密码标准。
DES使用64位的分组长度和密钥长度（其中有8位奇偶校验位），它的加密过程包括16轮的迭代，每轮使用一个48位的子密钥。解密过程与加密过程类似，只不过密钥的使用顺序是相反的，即从第16轮到第1轮。
\begin{figure}[ht]
	\centering
	\includegraphics[width=0.4\textwidth, keepaspectratio]{3-1.pdf}
	\caption{DES加密流程}
\end{figure}

\subsubsection{初始置换}
初始置换（Initial Permutation, IP） 是加密过程的第一步，它对明文块（64 位）按一个固定的矩阵进行比特位置重排，
$$
IP =
\begin{bmatrix}
58 & 50 & 42 & 34 & 26 & 18 & 10 & 2 \\
60 & 52 & 44 & 36 & 28 & 20 & 12 & 4 \\
62 & 54 & 46 & 38 & 30 & 22 & 14 & 6 \\
64 & 56 & 48 & 40 & 32 & 24 & 16 & 8 \\
57 & 49 & 41 & 33 & 25 & 17 & 9  & 1 \\
59 & 51 & 43 & 35 & 27 & 19 & 11 & 3 \\
61 & 53 & 45 & 37 & 29 & 21 & 13 & 5 \\
63 & 55 & 47 & 39 & 31 & 23 & 15 & 7 \\
\end{bmatrix}
\quad
IP^{-1} =
\begin{bmatrix}
40 & 8  & 48 & 16 & 56 & 24 & 64 & 32 \\
39 & 7  & 47 & 15 & 55 & 23 & 63 & 31 \\
38 & 6  & 46 & 14 & 54 & 22 & 62 & 30 \\
37 & 5  & 45 & 13 & 53 & 21 & 61 & 29 \\
36 & 4  & 44 & 12 & 52 & 20 & 60 & 28 \\
35 & 3  & 43 & 11 & 51 & 19 & 59 & 27 \\
34 & 2  & 42 & 10 & 50 & 18 & 58 & 26 \\
33 & 1  & 41 & 9  & 49 & 17 & 57 & 25 \\
\end{bmatrix}
$$

\subsubsection{轮函数}
轮函数（Round Function）是DES加密的核心部分，它将32位的输入数据与48位的子密钥进行处理，输出32位的结果。
轮函数的主要步骤包括：
\begin{enumerate}
    \item 扩展置换（Expansion Permutation），将32位输入扩展为48位。
    \item S盒替换（S-Box Substitution），将48位数据分为8组，每组6位，每组输出4位。
    \item P盒置换（P-Box Permutation），32位的置换。
\end{enumerate}
\begin{figure}[H]
	\centering
	\includegraphics[width=0.25\textwidth, keepaspectratio]{3-2.pdf}
	\caption{轮函数结构}
\end{figure}
$$
E =
\begin{bmatrix}
32 & 1  & 2  & 3  & 4  & 5  \\
4  & 5  & 6  & 7  & 8  & 9  \\
8  & 9  & 10 & 11 & 12 & 13 \\
12 & 13 & 14 & 15 & 16 & 17 \\
16 & 17 & 18 & 19 & 20 & 21 \\
20 & 21 & 22 & 23 & 24 & 25 \\
24 & 25 & 26 & 27 & 28 & 29 \\
28 & 29 & 30 & 31 & 32 & 1
\end{bmatrix}
\quad
P =
\begin{bmatrix}
16 & 7 & 20 & 21 \\
29 & 12 & 28 & 17 \\
1  & 15 & 23 & 26 \\
5  & 18 & 31 & 10 \\
2  & 8  & 24 & 14 \\
32 & 27 & 3  & 9 \\
19 & 13 & 30 & 6 \\
22 & 11 & 4  & 25
\end{bmatrix}
$$

S 盒中每一个元素都是 4 bit 的比特串，输入以 1, 6, 2, 3, 4, 5 的顺序重排序，然后输出表中对应位置的值。
$$
S1 =
\begin{bmatrix}
14 & 4 & 13 & 1 & 2 & 15 & 11 & 8 & 3 & 10 & 6 & 12 & 5 & 9 & 0 & 7 \\
0  & 15 & 7 & 4 & 14 & 2 & 13 & 1 & 10 & 6 & 12 & 11 & 9 & 5 & 3 & 8 \\
4  & 1  & 14 & 8 & 13 & 6 & 2 & 11 & 15 & 12 & 9 & 7 & 3 & 10 & 5 & 0 \\
15 & 12 & 8  & 2 & 4  & 9 & 1 & 7  & 5  & 11 & 3 & 14 & 10 & 0 & 6 & 13
\end{bmatrix}
$$

$$
S2 =
\begin{bmatrix}
15 & 1 & 8 & 14 & 6 & 11 & 3 & 4 & 9 & 7 & 2 & 13 & 12 & 0 & 5 & 10 \\
3  & 13 & 4 & 7  & 15 & 2  & 8 & 14 & 12 & 0 & 1 & 10 & 6 & 9 & 11 & 5 \\
0  & 14 & 7 & 11 & 10 & 4  & 13 & 1 & 5  & 8 & 12 & 6  & 9 & 3 & 2  & 15 \\
13 & 8  & 10 & 1 & 3  & 15 & 4 & 2 & 11 & 6 & 7 & 12 & 0 & 5 & 14 & 9
\end{bmatrix}
$$

$$
S3 =
\begin{bmatrix}
10 & 0 & 9 & 14 & 6 & 3 & 15 & 5 & 1 & 13 & 12 & 7 & 11 & 4 & 2 & 8 \\
13 & 7 & 0 & 9  & 3 & 4 & 6  & 10 & 2 & 8  & 5  & 14 & 12 & 11 & 15 & 1 \\
13 & 6 & 4 & 9  & 8 & 15 & 3 & 0 & 11 & 1  & 2  & 12 & 5 & 10 & 14 & 7 \\
1  & 10 & 13 & 0 & 6 & 9  & 8 & 7 & 4  & 15 & 14 & 3  & 11 & 5 & 2  & 12
\end{bmatrix}
$$

$$
S4 =
\begin{bmatrix}
7 & 13 & 14 & 3 & 0 & 6 & 9 & 10 & 1 & 2 & 8 & 5 & 11 & 12 & 4 & 15 \\
13 & 8 & 11 & 5 & 6 & 15 & 0 & 3 & 4 & 7 & 2 & 12 & 1 & 10 & 14 & 9 \\
10 & 6 & 9 & 0 & 12 & 11 & 7 & 13 & 15 & 1 & 3 & 14 & 5 & 2 & 8 & 4 \\
3 & 15 & 0 & 6 & 10 & 1 & 13 & 8 & 9 & 4 & 5 & 11 & 12 & 7 & 2 & 14
\end{bmatrix}
$$

$$
S5 =
\begin{bmatrix}
2 & 12 & 4 & 1 & 7 & 10 & 11 & 6 & 8 & 5 & 3 & 15 & 13 & 0 & 14 & 9 \\
14 & 11 & 2 & 12 & 4 & 7 & 13 & 1 & 5 & 0 & 15 & 10 & 3 & 9 & 8 & 6 \\
4 & 2 & 1 & 11 & 10 & 13 & 7 & 8 & 15 & 9 & 12 & 5 & 6 & 3 & 0 & 14 \\
11 & 8 & 12 & 7 & 1 & 14 & 2 & 13 & 6 & 15 & 0 & 9 & 10 & 4 & 5 & 3
\end{bmatrix}
$$

$$
S6 =
\begin{bmatrix}
12 & 1 & 10 & 15 & 9 & 2 & 6 & 8 & 0 & 13 & 3 & 4 & 14 & 7 & 5 & 11 \\
10 & 15 & 4 & 2 & 7 & 12 & 9 & 5 & 6 & 1 & 13 & 14 & 0 & 11 & 3 & 8 \\
9 & 14 & 15 & 5 & 2 & 8 & 12 & 3 & 7 & 0 & 4 & 10 & 1 & 13 & 11 & 6 \\
4 & 3 & 2 & 12 & 9 & 5 & 15 & 10 & 11 & 14 & 1 & 7 & 6 & 0 & 8 & 13
\end{bmatrix}
$$

$$
S7 =
\begin{bmatrix}
4 & 11 & 2 & 14 & 15 & 0 & 8 & 13 & 3 & 12 & 9 & 7 & 5 & 10 & 6 & 1 \\
13 & 0 & 11 & 7 & 4 & 9 & 1 & 10 & 14 & 3 & 5 & 12 & 2 & 15 & 8 & 6 \\
1 & 4 & 11 & 13 & 12 & 3 & 7 & 14 & 10 & 15 & 6 & 8 & 0 & 5 & 9 & 2 \\
6 & 11 & 13 & 8 & 1 & 4 & 10 & 7 & 9 & 5 & 0 & 15 & 14 & 2 & 3 & 12
\end{bmatrix}
$$

$$
S8 =
\begin{bmatrix}
13 & 2 & 8 & 4 & 6 & 15 & 11 & 1 & 10 & 9 & 3 & 14 & 5 & 0 & 12 & 7 \\
1 & 15 & 13 & 8 & 10 & 3 & 7 & 4 & 12 & 5 & 6 & 11 & 0 & 14 & 9 & 2 \\
7 & 11 & 4 & 1 & 9 & 12 & 14 & 2 & 0 & 6 & 10 & 13 & 15 & 3 & 5 & 8 \\
2 & 1 & 14 & 7 & 4 & 10 & 8 & 13 & 15 & 12 & 9 & 0 & 3 & 5 & 6 & 11
\end{bmatrix}
$$

\subsubsection{密钥扩展}
密钥扩展（Key Schedule）是DES加密的一个重要步骤，它将64位的密钥扩展为16个48位的子密钥。
密钥扩展的过程包括：
\begin{enumerate}
    \item PC-1（Permuted Choice 1），将64位密钥转换为56位。
    \item 分为两部分（C和D），每部分28位，进行16轮迭代。
    \item 两部分分别循环左移（Left Shift），当i=1，2，9，16时，左移一次，其他轮次左移两次。
    \item 使用PC-2生成48位子密钥。
\end{enumerate}

\begin{figure}[H]
	\centering
	\includegraphics[width=0.4\textwidth, keepaspectratio]{3-3.pdf}
	\caption{密钥扩展}
\end{figure}

$$
PC_1 =
\begin{bmatrix}
57 & 49 & 41 & 33 & 25 & 17 & 9 \\
1  & 58 & 50 & 42 & 34 & 26 & 18 \\
10 & 2  & 59 & 51 & 43 & 35 & 27 \\
19 & 11 & 3  & 60 & 52 & 44 & 36 \\
63 & 55 & 47 & 39 & 31 & 23 & 15 \\
7  & 62 & 54 & 46 & 38 & 30 & 22 \\
14 & 6  & 61 & 53 & 45 & 37 & 29 \\
21 & 13 & 5  & 28 & 20 & 12 & 4
\end{bmatrix}
\quad
PC_2 =
\begin{bmatrix}
14 & 17 & 11 & 24 & 1  & 5  \\
3  & 28 & 15 & 6  & 21 & 10 \\
23 & 19 & 12 & 4  & 26 & 8  \\
16 & 7  & 27 & 20 & 13 & 2  \\
41 & 52 & 31 & 37 & 47 & 55 \\
30 & 40 & 51 & 45 & 33 & 48 \\
44 & 49 & 39 & 56 & 34 & 53 \\
46 & 42 & 50 & 36 & 29 & 32
\end{bmatrix}
$$

\subsubsection{TDES}
DES极大缺陷是它使用了短密钥，为了克服这一点，Tuchman在1979年提出了三重DES（Triple DES, TDES），
它通过对数据进行三次DES加密来提高安全性，密钥长度为168位，当$k_1=k_2=k_3$，TDES退化成DES，
$k=(k_1, k_2, k_3)$，加密过程为：
$$ TDES(m)=DES_{k_3}(DES^{-1}_{k_2}(DES_{k_1}(m))) $$